In our earlier work \cite{DTS-2025} we addressed the question of whether
a set $\mathcal{S}$ of normalised (sequential) scenarios, whose
sequences of events are different permutations of the same set, can be
realised as a single DTS that is
semantically equivalent to $\mathcal{S}$, i.e., describes the same set of
behaviours as $S$. This can be done only if there is a distance
relation whose set of projections is the same as $\mathcal{S}$. We
showed how to obtain such a distance relation and a corresponding
DTS \cite{DTS-2025}.
%(see \Sec{sec-dts} for a brief summary).

In this section we tackle the problem of realisability of
scenarios \emph{as expressions}: given a set, $\mathcal{S}$,
of scenarios, we want to find an expression that realises
$\mathcal{S}$.
Despite similarities, there is a fundamental difference between
this problem and the one previously solved \cite{DTS-2025}: the
members of $\mathcal{S}$ could be projections that are obtained from
extended distance relations of different distributed
scenarios, in which case $\mathcal{S}$ cannot be realised by a single
DTS, but by a choice expression that involves more than one DTS. The
challenge is to divide $\mathcal{S}$ in such a way that each
partition could be realised by a DTS, while trying to avoid the trivial
solution in which every member of S is a separate alternative.
More precisely:
\vspace{-0.5em}
\begin{quote}
Given a finite set $\mathcal{S}=\{\SeqS_1,\dots,\SeqS_n\}$ of
scenarios of \emph{possibly different lengths} over
$\Sigma$, such that
(i) $\ESeq{\SeqS_i}\neq \ESeq{\SeqS_j}$ (for $1\leq
i,j\leq n$ and $i\neq j$)\footnote{In \Sec{sec-relax} we
show how to relax the assumption.}, and (ii)
each event of $\SeqS_i$ ($1\leq i\leq n$) has a
unique occurrence in $\SeqS_i$,
we want to obtain an expression $E$
over $\Sigma$, such that $\Sem{E} = \Sem{\mathcal{S}}=\bigcup_{1\leq
i \leq n} \Sem{\SeqS_i}$.
In general, $\mathcal{S}$ can be realised by more than one
expression.
Our goal is to find an expression with a
reasonably small number of alternatives, while decreasing the number
of trivial alternatives.
\end{quote}
\vspace{-0.6em}
We say $\mathcal{S}$ is \emph{realisable by} $E$, or
$E$ is a \emph{realisation} of $\mathcal{S}$, if $\mathcal{S}$ is
realised by $E$.
\vspace{-0.2em}
\begin{example}
\label{ex-one-flag}
Given the set, $\mathcal{S}=\{\SeqS_1,\SeqS_2,\SeqS_3,\SeqS_4\}$, of
scenarios, where

\noindent
$\SeqS_1=(\mathtt{start}~ \mathtt{finish_1}~ \mathtt{finish_2}~ \mathtt{flag_1},\emptyset)$,
$\SeqS_2=(\mathtt{start}~ \mathtt{finish_1}~  \mathtt{flag_1}~ \mathtt{finish_2},\emptyset)$,
$\SeqS_3=(\mathtt{start}~ \mathtt{finish_2}~ \mathtt{finish_1}~ \mathtt{flag_2},\emptyset)$,
%and
$\SeqS_4=(\mathtt{start}~ \mathtt{finish_2}~ \mathtt{flag_2}~ \mathtt{finish_1},\emptyset)$,
there is no DTS that realises $\mathcal{S}$.
But $\mathcal{S}$ can be divided into two sets
$\mathcal{S}_{1}=\{\SeqS_1,\SeqS_2\}$ and
$\mathcal{S}_2=\{\SeqS_3,\SeqS_4\}$ based on the sets of events of its
members. The set ${\mathcal{S}}_{1}$ can be realised by DTS
$E=\D(\xi_1,\xi_2)=\xi_1||\xi_2$, where
$\xi_1=(\mathtt{start}~ \mathtt{finish_1}~ \mathtt{flag_1},\emptyset)$
and
$\xi_2=(\mathtt{start}~ \mathtt{finish_1}~ \mathtt{finish_2},\emptyset)$.
Similarly, $\mathcal{S}_2$ is realised by DTS
$E'=\xi_3||\xi_4$, where
$\xi_3=(\mathtt{start}~ \mathtt{finish_2}~ \mathtt{flag_2},\emptyset)$
and
$\xi_4=(\mathtt{start}~ \mathtt{finish_2}~ \mathtt{finish_1},\emptyset)$.
So $\mathcal{S}$ is realised by the expression
$E\BAR E'=(\xi_1||\xi_2)|(\xi_3||\xi_4)$.
\end{example}
\vspace{-0.5em}
%------------------------
\begin{definition}
  \label{def-DR-S}
Let $\mathcal{S}=\{\SeqS_1,\dots,\SeqS_n\}$ be a finite set of
normalised scenarios
%over $\Sigma_{\mathcal{S}}$
and
$\{\DR{\SeqS_1},\dots,\DR{\SeqS_n}\}$ be the set of
corresponding stable distance relations, where
$\DR{\SeqS_i}=(\PR{\SeqS_i},\DF{\SeqS_i})$ ($1\leq i\leq n$).
\emph{The distance relation of} $\mathcal{S}$, $\DR{\mathcal{S}}$, is
the pair $(\ParS{\mathcal{S}},\DF{\mathcal{S}})$, where
$\ParS{\mathcal{S}}=\bigcap_{1\leq i\leq n} \PR{\SeqS_i}$ and,
for every $(e,e')$ in $\ParS{\mathcal{S}}$,
$\DF{\mathcal{S}}(e,e')$ is
the smallest interval that includes each $\DF{\SeqS_i}(e,e')$.
%$\DR{\mathcal{S}}$ is stable and cannot be extended.
\end{definition}
\vspace{-0.3em}
\noindent
The first step\label{initial-partitioning-step}
 in solving the problem posed at the beginning of this
section is to divide
$\mathcal{S}$ in such a way that all the members of each partition are
scenarios of equal lengths over the same set of events.
Then we solve the problem for each partition, independently.
Our general strategy is the following: whenever a partition is not
realisable by a single DTS, we divide it
%If a partition is not realisable by a single DTS, we divide it
repeatedly, so that at the end each partition will be realised by a
DTS and will correspond to one alternative in the overall choice expression.
This process is guided by the results of our analysis based on
$\PR{\mathcal{S}}$ and
%the constraints of $\mathcal{S}$,
$\DF{\mathcal{S}}$,
described in sections \ref{sec-no-time}
and \ref{sec-full}. Additionally, some divisions are simply necessary, as
described below.
%But first, an important definition:
%-------------
\vspace{-1em}
\subsubsection{Gaps}\label{paragraph-gaps}
Let $S_{(e, e')}=\{s_i\in \mathcal{S}|e\PR{\SeqS_i}e'\}$.
Let $s_1,\dots,s_m$ be a sequence of members of $S_{(e, e')}$ such
that, for every $1\leq i\leq m$, $[l_i, h_i]=\DF{s_i}(e, e’)$
and $l_i \leq l_{i+1}$.
Let $H_{r}=\max \{h_j|1\leq j \leq r\}$.
If there is a $k$, $1 \leq k < m$, such that
$H_{k} < l_{k+1}$, $\DF{S}(e, e’)$ contains a gap.
Then $\mathcal{S}$ is said to be \emph{with gaps}.
An $\mathcal{S}$ with gaps cannot be realised by a single DTS:
this is because the interval between $e$ and $e'$ in
$\DF{\mathcal{S}}$ would be
$[l_1,H_{m}]$, which would allow also behaviours
 in which the time distance
 between $e$ and $e'$ is between $H_k$ and $l_{k+1}$. Such behavious are
  allowed by none of the members of $\mathcal{S}$.
  So $\mathcal{S}$ must be divided into partitions with no gaps.
%, but would be allowed by any expression that realises $\mathcal{S}$.

Whenever a partition is divided into finer partitions (see \Sec{sec-no-time}
 and \Sec{sec-full}), we must check for the existence of gaps in any
 newly formed partition.%  (see \Sec{sec-algo}).% and divide accordingly.


% In general, it might be necessary to divide a partition further into
% subgroups (see \Sec{sec-no-time}
% and \Sec{sec-full}). So we must check for the existence of gaps in
% any newly formed partition and divide accordingly.

In the rest of the paper we assume that all scenarios in
$\mathcal{S}$ are over the same set, $\Sigma_{\mathcal{S}}$, of
events, and that each $\SeqS\in\mathcal{S}$
contains all the events from $\Sigma_{\mathcal{S}}$.
Moreover, we assume that any divisions necessitated by
gaps have already been carried out, so that $\mathcal{S}$ is without gaps.
%, and moreover the intervals of $\DF{\mathcal{S}}$ are without gaps.
We describe our method in two steps:
%in the first step
first, we tackle the
problem \emph{without accounting for constraints}.
Because constraints will not play a role,
in our examples we will discuss scenarios with only default constraints.
%we assume each member of $\mathcal{S}$ is of the form
%$(\E,\emptyset)$: it has only default constraints.
%---------------
